\section{Seven dimensions from one enlarged line
class}

The
\href{https://www.math.rwth-aachen.de/~Gabriele.Nebe/LATTICES/P48p.html}{catalogue
lattice \(P_{48p}\)} is even unimodular of dimension 48 and minimum
squared norm six. Its theta series determines a minimal shell of
52,416,000 vectors~\cite{p48p-catalogue,nebe}.

\subsection{Disjoint blocks from overlapping images}

The new finite object is a class of 7,077 minimal lines. Integer
representatives \(c\) in the catalogue Gram basis satisfy \(cGc^t=6\)
and

\[
|cGc'^t|\le2\qquad(c\ne\pm c').
\]

Exact lattice isometries give further classes. Select a family of images
and assign every line to its first occurrence. This produces disjoint
line subsets with union size \(S\). Taking both signs of every line
produces compatible head blocks.

Use the root codes \(A_1,A_2,D_3,D_4,D_5,E_6,E_7\) as tails in
dimensions one through seven. Their sizes are \(2,6,12,24,40,72,126\).
Assign each head block to a different antipodal tail line. A used mother
line loses its two equatorial points and contributes four cap points at
head squared norm \(3/4\) and tail squared norm \(1/4\). The count is

\[
K(48+k)\ge52416000-2S+4S+\tau_k=52416000+2S+\tau_k.
\]

\begin{longtable}[]{@{}rrrr@{}}
\toprule
\(k\) & Selected images & Union size \(S\) & Tail points
\(\tau_k\)\tabularnewline
\midrule
\endhead
1 & 1 & 7,077 & 2\tabularnewline
2 & 3 & 21,231 & 6\tabularnewline
3 & 6 & 42,459 & 12\tabularnewline
4 & 12 & 84,874 & 24\tabularnewline
5 & 20 & 141,328 & 40\tabularnewline
6 & 36 & 253,851 & 72\tabularnewline
7 & 63 & 442,507 & 126\tabularnewline
\bottomrule
\end{longtable}

\subsection{Why the lift preserves all angular constraints}

Let $x$ be a selected mother vector divided by $\sqrt6$, and let $p_i$
represent its assigned unit tail line. The four points replacing
$\{\pm x\}$ are
\[
 \left(\varepsilon\frac{\sqrt3}{2}x,
       \delta\frac12p_i\right),\qquad \varepsilon,\delta\in\{\pm1\}.
\]
They have unit norm. Within a block, distinct head lines have absolute
product at most $1/3$; in the entire minimal shell the corresponding
bound is $1/2$. These two facts give all cap inequalities:
\[
\begin{array}{ll}
\text{different lines in one block:}&(3/4)(1/3)+1/4=1/2,\\
\text{same head, opposite tails:}&3/4-1/4=1/2,\\
\text{opposite heads:}&-3/4+1/4=-1/2,\\
\text{different blocks:}&(3/4)(1/2)+(1/4)(1/2)=1/2.
\end{array}
\]
The last line uses distinct tail lines. Reusing a tail line for arbitrary
different blocks would permit $5/8$. A cap meets a retained equatorial
point with product at most $\sqrt3/4$, and a pure tail with product at
most $1/2$. The retained equator and pure tails already satisfy the
kissing condition. Thus the construction and its count are established
by the class inequalities and the lattice minimum.

\subsection{Separating the two sources of gain}

The gain over the public 7,069-line class~\cite{kissingnumbers} has two
sources: eight additional mother lines and improved choices of isometric
images. Apply the selected images to the original class and call the
union size $S_0$; let $S_{\rm ref}$ be the comparison construction's
union size. Then
\[
 N-N_{\rm ref}=2(S-S_0)+2(S_0-S_{\rm ref}).
\]
The first term holds the images fixed and measures class growth. The
second holds the original class fixed and measures image selection.

\begin{center}
\begin{tabular}{rrrr}
\toprule
Dimension&Class growth&Image selection&Total increase\\\midrule
49&16&0&16\\
50&48&2&50\\
51&96&18&114\\
52&192&64&256\\
53&320&154&474\\
54&572&296&868\\
55&996&414&1410\\\bottomrule
\end{tabular}
\end{center}

The three images in dimension 50 are disjoint. In dimension 55,
overlap reduces the $8\cdot63$ possible new lines to 498 distinct
lines; image selection on the original class supplies another 207.
This is why enlarging one class and maximizing the union are
complementary parts of the construction.
